\subsection{Genuine grading of a homogeneous ideal quotient}
\label{subsec:homogeneous-quotient-grading}

\begin{theorem}[Actual quotient grading, with every original homogeneous lift]
\label{graded:homogeneous-ideal-quotient-grading}
\lean{EnvelopeProperty.GradedQuotient.grading,
EnvelopeProperty.GradedQuotient.componentMap,
EnvelopeProperty.GradedQuotient.componentMap_coe,
EnvelopeProperty.GradedQuotient.componentMap_surjective,
EnvelopeProperty.GradedQuotient.instGradedMonoid,
EnvelopeProperty.GradedQuotient.isInternal,
EnvelopeProperty.GradedQuotient.gradedRing,
EnvelopeProperty.GradedQuotient.gradedMk,
EnvelopeProperty.GradedQuotient.gradedMk_apply,
EnvelopeProperty.GradedQuotient.instGradedRingHomogeneous,
EnvelopeProperty.GradedQuotient.irrelevant_le_map,
EnvelopeProperty.GradedQuotient.map_irrelevant}
\leanok
Let $A$ be a commutative algebra over an arbitrary commutative coefficient
ring $R$, internally graded by an additive index monoid. Let $I$ be an
actual homogeneous ideal, without reducedness or primeness assumptions.
The actual quotient ring $A/I$ has its genuine internal grading given
by the images of the original homogeneous components. The original
quotient ring homomorphism is genuinely graded and surjective on every
homogeneous component. Every homogeneous quotient element therefore
has an actual lift of the same degree. This construction preserves the
original ideal, including all its nilpotents. For a natural-number
grading its actual irrelevant ideal is exactly the image of the original
irrelevant ideal under this genuine quotient map.
\end{theorem}
\begin{proof}
Define each quotient component as the actual submodule image under the
original quotient algebra homomorphism. The quotient of the original
homogeneous one lies in degree zero, and the actual images of products
of homogeneous elements have the required sum degree. This proves the
actual graded multiplicative structure on these component images.

The quotient homomorphism gives a surjective additive map on each
component by the definition of its genuine submodule image. The actual
componentwise direct-sum map is therefore surjective. The canonical
recomposition homomorphisms commute with this entire direct-sum map
and the original quotient map, as follows on each actual homogeneous
summand and hence on every finite direct sum.

Every quotient element is represented by an original element. Decompose
that representative in the genuine original internal grading and map
its components to the quotient. The recomposition identity proves
surjectivity of the canonical quotient-component recomposition.

For injectivity, take a quotient direct sum whose recomposition is
zero and lift its finitely many homogeneous components to an actual
original direct sum using the proved componentwise surjectivity. Its
original recomposition lies in the actual ideal $I$, by the quotient
kernel criterion. Homogeneity of $I$ puts every original homogeneous
component in $I$. The original internal decomposition identifies these
components with the selected original direct-sum entries. Their quotient
images are all zero, so the whole quotient direct sum is zero. Applying
this to a difference proves genuine injectivity.

Thus the actual canonical quotient-component recomposition is
bijective. Its true inverse supplies an internal decomposition,
combined with the preceding actual graded multiplication. The original
quotient map preserves every homogeneous component by construction,
and the same actual submodule image witnesses give the stated lifts.
Neither an internal decomposition nor a graded quotient certificate is
assumed; no radical replacement, domain, field, noetherian or finite
generation condition is used. For the natural-number grading, each
positive quotient component lifts in the same positive degree, so it
lies in the actual image of the original irrelevant ideal. Conversely
the true graded quotient map takes every positive original component
into the quotient irrelevant ideal. The actual positive-component
generator criterion for irrelevant ideals then proves the asserted
entire ideal-image equality.
\end{proof}
